% TLDR: State the judge-only contribution and distinguish completed evidence from pending collection.
\section{Discussion and Current Status}
The study asks which training objective yields the best practical judge of human ratings under matched resources.
Proper scoring rules identify a distributional population target, but finite data, model restrictions and optimization can alter the empirical comparison.
The historical mixture analysis leaves the RPS-versus-Pearson advantage unresolved.
The selected scalar judge also shows why strong rank association can coexist with systematic leniency.
Neither result establishes a winner for a future native-probability judge.

The three workstreams have different roles: historical ratings support retrospective analyses, the ongoing packet prepares new human evidence, and size planning determines what additional evidence is feasible.
The replacement packet plans 85 artifacts and 170 independent ratings; its completed manifest and the prospective split amendment remain unfinished.
Exact judge pins, model-repeat counts, practical success thresholds and the sparse-assignment uncertainty procedure must be fixed before test inspection.
Warm-up responses are collection progress, not completed main-round evaluation.

A successful study would support a bounded claim about predicting ratings under specified inputs, assistance conditions and rater recruitment.
It would not equate agreement with semantic truth, establish every artifact's correctness, or guarantee transfer to unrepresented tasks, solvers or people.
Independent semantic controls and checked proof evidence address complementary questions.
The paper therefore contributes a defined comparison and traceable historical evidence while reserving fresh empirical conclusions for the new collection.
